% Separate post-CP27 candidate. The nonsquare whole-formal branch remains open.
\subsection{The finite-character gate at the first wild cyclotomic index}

\begin{lemma}[An unweighted iterated product-height-two bound]
\label{lem:B-RD77-kernel}
\status{Candidate; complete proof supplied, independent audits pending.}
Let \(R\) be a complete discrete valuation ring of mixed
characteristic \((0,p)\), with algebraically closed residue
field and \(p\ge5\). Write \(a=v_R(p)\). Let \(f:F\to G\)
be a homomorphism of finite-height formal groups with no nonzero
finite flat subgroup in its kernel and with finite special-fibre
kernel. Assume that the actual special fibre \(G_k\) is a product
of one-dimensional height-two formal groups. Define
\[
 q_0=1,\qquad q_{j+1}=\min\{a+q_j,1+pq_j\}.
                                                        \tag{B-RD77.1}
\]
If \(q_u\ge a/(p-1)\), then \(F^u\ker f_k=0\).
\end{lemma}
\begin{proof}
The invariant-differential identity
\(U([p]X)J_{[p]}(X)=pU(X)\) shows that every monomial with
an exponent not divisible by \(p\) has coefficient divisible
by \(p\). The product-height-two special fibre makes every
coefficient of degree less than \(p^2\) vanish modulo the
uniformizer, except that the linear term is already \(pX\).
For an augmentation vector of coefficient valuation at least
\(q\ge1\), this gives the lower bound
\[
 \min\{a+q,1+pq,p^2q\}
       =\min\{a+q,1+pq\}.                             \tag{B-RD77.2}
\]
This is the unweighted counterpart of (B-RD64.5); it concerns
the actual target and is invariant under integral coordinate
changes. Iteration gives \(q_j\).

For precision, the normalized logarithm coefficient bound is
\(v(l_\alpha)\ge-a\lfloor\log_p|\alpha|\rfloor\).
Therefore, at the stated closed threshold, coefficientwise
substitution gives
\[
       p^u\log_G(\pi h)\in\pi^{q_u}R[[t]]^{\dim G}
       \quad(h\in tR[[t]]^{\dim G}).                  \tag{B-RD77.3}
\]
The exact integral Cartier action, bounded projection and
saturation argument in Lemma~\ref{lem:B-RD64-weighted-kernel}
uses only this last inclusion after its initial coefficient
calculation. It applies unchanged: for
\(J=\pi^{\lceil a/(p-1)\rceil}R\), the integral
\(\Omega\)-action preserves \(J\), the bounded projection
of \(\mathbf f^u z\) lies in \(J[[\Delta]]\), and the
absence of a nonzero finite flat kernel identifies the source
Cartier module with the target intersection. Thus the saturation
quotient is killed by \(\mathbf f^u\). The cited identification
of that quotient with the Dieudonne module of the special kernel
makes \(\mathbf f\) relative Frobenius. These are precisely
\cite[Propositions~9--12 and Lemma~3, pp.~674--676;
Proposition~16 and proof, p.~681; Theorem~2(8), p.~668]
{Bondarko2006Generic}, with the closed coefficientwise boundary
already proved in Lemma~\ref{lem:B-RD50-boundary}.
No target thickness is inferred from a source or an isogeny.
\end{proof}

\begin{proposition}[Square finite determinants at the first wild index]
\label{prop:B-RD77-square}
\status{Candidate; complete proof supplied, independent audits pending.}
Let \(p\ge5\), and let \(K\) be complete discretely valued
of mixed characteristic \((0,p)\), with algebraically closed
residue field and \(e=p(p-1)\). Let \(M,N\) be whole
potentially supersingular height-two groups with common
determinant \(\chi_p\psi\). If the finite
\(\mu_{p-1}\)-valued character \(\psi\) is a square in
\(\operatorname{Hom}_{\rm cont}(I_K,\mu_{p-1})\), then
\[
             M[p^3]\simeq N[p^3]
   \quad\Longrightarrow\quad
   \text{their finite special inertia characters agree}.
                                                        \tag{B-RD77.4}
\]
\end{proposition}
\begin{proof}
Twist both groups by the inverse of a square root of \(\psi\).
This is the actual finite twist of
Proposition~\ref{prop:B-RD50-square}; it preserves the given
torsion comparison and preserves inequality of the two special
characters. Their finite special determinants are now one.
The finite height-two inertia groups are cyclic subgroups of
\(\mathbf F_{p^2}^{\times}\), and their norm-one condition
makes their common faithful good degree \(m\) divide \(p+1\).
Thus, over that actual good field,
\[
             a=em=p(p-1)m\le p(p^2-1).               \tag{B-RD77.5}
\]

Suppose the special characters differ. Their norm-one
Frobenius eigenvalue orbits are disjoint, so an integral tame
polynomial projector separates the two special factors, just
as in B50. Let \(H\) be the finite-flat closure of the given
generic graph in \(\mathcal M[p^3]\times\mathcal N[p^3]\).
The perfect generic alternating pairings, and their inverse
duality maps, extend to perfect integral BT pairings by
\cite[Theorem~4 and Corollaries~1--2, p.~181]{Tate1967PDivisible}.
After scaling one by the unit multiplier of the torsion map,
the graph closure is self-orthogonal, including on the special
fibre. The integral projector splits \(H_k\) as a product.
Uniqueness of a rank-\(p^r\) subgroup in each one-dimensional
connected formal group and self-orthogonality force
\[
            H_k=\ker F_{M_k}^{3}\times\ker F_{N_k}^{3}.
                                                        \tag{B-RD77.6}
\]

The actual quotient \(Q=(\mathcal M\times\mathcal N)/H\)
has special fibre \(M_k^{(p^3)}\times N_k^{(p^3)}\).
The map
\[
 \mathcal M\times\mathcal N\longrightarrow Q\times\mathcal M,
 \qquad z\longmapsto(qz,\operatorname{pr}_{\mathcal M}z)
                                                        \tag{B-RD77.7}
\]
has no nonzero finite flat kernel, since its generic kernel is
zero, and has special kernel \(0\times\ker F_{N_k}^{3}\).
Its actual target is a product of three height-two special
groups. Apply the preceding lemma with \(u=2\). Here
\[
 q_1=p+1,\qquad
 q_2=\min\{a+p+1,p^2+p+1\}\ge a/(p-1),              \tag{B-RD77.8}
\]
because \(a\ge p(p-1)\), \(a\le p(p^2-1)\), and both
entries of the minimum are at least that threshold. The
resulting assertion that \(F^2\) kills \(\ker F^3\) is
impossible: its image has rank \(p\). This proves the proposition.
\end{proof}

\begin{corollary}[Two remaining patterns at \(p=5,e=20\)]
\label{cor:B-RD77}
\status{Candidate; complete proof supplied, independent audits pending.}
In the full B.4 height-two setting with perfect residue field,
\(p\ge5\), \(\mu_p\subset K\), and \(e=p(p-1)\),
\[
                    2\le\nu_{\min}(p,p(p-1),2)\le3.
                                                        \tag{B-RD77.9}
\]
Any actual \(p^3\)-torsion failure must compare two whole
potentially supersingular groups with a nonsquare common finite
determinant. After a simultaneous finite twist, its value may
be normalized to be primitive of order \(p-1\).

For \(p=5,e=20\), the joint faithful tame degree of such a
remaining pair is \(24\). If \(\zeta\) is a primitive
\(24\)th root whose sixth power is the prescribed primitive
norm, the special characters have distinct types among
\[
 A=\{\zeta,\zeta^5\},\qquad
 B=\{\zeta^9,\zeta^{21}\},\qquad
 C=\{\zeta^{13},\zeta^{17}\}.                         \tag{B-RD77.10}
\]
Up to interchange and a simultaneous quadratic twist, the only
remaining type pairs are \(\{A,B\}\) and \(\{A,C\}\).
This is a necessary finite-character list, not an existence
assertion for torsion-congruent groups.
\end{corollary}
\begin{proof}
The B63 actual \(p^2\)-torsion construction applies because
\(p(p-1)>p\), and gives the lower bound two. At actual level
\(p^3\), Lemma~\ref{lem:B-RD76-determinant} applies because
\(v_p(e)=1\), and Corollary~\ref{cor:B-RD76} removes every
case except the whole-formal comparison. The preceding proposition
removes its square finite determinant.

For the upper bound, use actual \(p^4\)-torsion. Dimension
detection again applies, while the mixed comparison would require
\(e>p^{2\cdot4-3-1}=p^4\), contrary to \(e=p(p-1)\).
For two whole formal groups, the common determinant is fixed.
In Proposition~\ref{prop:B-RD64-formal} take
\(a=e(p^2-1)=p(p-1)(p^2-1)\). Its actual-target weighted
recurrence gives
\[
 q_1=2p-1,\qquad q_2=2p^2-1,\qquad q_3=2p^3-1
                  >p(p^2-1)=a/(p-1).                \tag{B-RD77.11}
\]
Thus level \(p^4\), or manuscript index three, suffices.

The normalization and finite-character enumeration do not
depend on the absolute ramification. The procyclic tame
character argument of Corollary~\ref{cor:B-RD50-remaining}
changes any nonsquare norm to a primitive one. The arithmetic
proof of Lemma~\ref{lem:B-RD51-orbits} then applies unchanged.
For \(p=5\), it permits joint degree \(8\) or \(24\), but
the primitive norm fibre in \(\mu_8\) has only one Frobenius
orbit, so different special characters require \(24\).
The three orbits and their simultaneous quadratic-twist action
are exactly the elementary computation in
Corollary~\ref{cor:B-RD51-p5}, giving (B-RD77.10) and the two
stated patterns. No earlier least-index statement at \(e=p-1\)
is imported to this wild index.
\end{proof}

\begin{remark}[The remaining wild comparison is still actual]
\label{rem:B-RD77-scope}
\status{Candidate scope; the two normalized patterns remain unresolved.}
The square exclusion uses only the actual quotient special fibre
and the coefficientwise logarithm bound; it does not transport a
source Hasse thickness through an isogeny. The remaining
\(p=5,e=20\) question is whether either of the displayed
nonsquare patterns occurs for an isomorphism of actual
\(125\)-torsion. The finite-character list and the upper index
three do not resolve it or supply a counterexample.

The retained exact supplement is
\path{code/reduction_descent/verify_first_wild_finite_patterns.py}.
Its first author run passed in \(0.6461968240037095\) seconds
at 23:15:03 UTC on 19 November 2026. For the \(24\) primes
from \(5\) through \(101\), it checks \(174\) norm-one
good-degree rows, \(348{,}000\) closed coefficientwise logarithm
inequalities, \(2{,}560{,}584\) low-degree coefficient bounds,
and all \(24\) weighted three-step upper bounds. The explicit
\(p=5\) primitive-norm fibres and quadratic-twist pair orbits
are retained separately. These finite arithmetic checks do
not imply existence of a BT group or a torsion congruence.
The complete code, output, finite ranges and pre-certificate
proof are retained; no executed computation failed.
\end{remark}
